\section{Pass order and protocol status}\label{si:passes}

The study is a fixed computational protocol with a recorded execution order, not a prospective registration: no dated pre-scoring protocol file exists for the empirical passes. The order below is established from the analysis code and its archived metadata. Passes 1 and 2 are labelled in the driver script \texttt{run\_ladder.py} and distinguished in \texttt{results/meta.json} by the keys \texttt{catch\_pass1} and \texttt{catch\_pass2}; passes 3 and 4 are separate scripts that import from that driver, which fixes their order after it.

\begin{table}[!htb]
\centering\small
\caption{Order in which the passes were run.}
\begin{tabularx}{\textwidth}{@{}clLL@{}}
\toprule
\# & Pass & Data treatment & Script / archived output \\
\midrule
1 & Coarse-regime catch & 240/120/5\,kt (through 1991 / 1992 / from 1993) & \texttt{run\_ladder.py} (pass 1); \texttt{fixed\_window\_scores.csv}, \texttt{rolling\_summary.csv} (rows \texttt{catch = regime}), \texttt{rolling\_forecasts\_regime.csv} \\
2 & Annual landings & Schijns et al.\ (2021) Table 1, t/1000 & \texttt{run\_ladder.py} (pass 2); same files, rows \texttt{catch = annual}; \texttt{rolling\_forecasts.csv} \\
2b & Survey start & Forecast initiated at the survey-scaled state, annual catch & \texttt{run\_survey\_start()} within pass 2; \texttt{survey\_start\_forecasts.csv} \\
3 & Capelin variants & Two-regime $r$ at the 1991 break; acoustic-index multiplier on observed years & \texttt{run\_capelin\_regime.py}, \texttt{run\_capelin\_index.py}; \texttt{capelin\_*\_summary.csv} \\
4 & Specification B & The ladder on the xteNCAM series alone, not pooled & \texttt{run\_xte.py}; \texttt{xte\_rolling\_summary.csv}, \texttt{xte\_fixed\_window\_scores.csv} \\
5 & Uncertainty layer & DM statistics and block bootstrap on the archived per-origin files & \texttt{campaign\_e1\_dm\_uncertainty\_regime.py}; \texttt{e1\_dm\_uncertainty\_regime.csv} \\
6 & Simulation & D1--D5 per \texttt{SPECIFICATION\_v4.md} (locked 2026-09-10 before the code was written); D6--D7 per Amendment 1 (locked 2026-09-11) & \texttt{tools/sim\_retention\_power.py}; \texttt{sim\_retention\_power.csv}, \texttt{sim\_misspecified\_D6D7.csv} \\
7 & This revision & Re-execution of passes 1--6 in a third environment; paired bootstrap confirmation; period decomposition; rolling parameter record; catch-timing sensitivity & \texttt{paperE1\_cod\_forecast\_ladder\_v69\_src/work/} in the repository \\
\bottomrule
\end{tabularx}
\end{table}

Three points bear on how the scores should be read. The five ladder modules were fixed before any pass was scored: the \texttt{SPECS} list in \texttt{run\_ladder.py} defines M1, M1b, M2, M3 and M4 once, Specification B imports that list rather than redefining it, and none was altered after a score was seen. The scoring core, the fixed windows, the rolling-origin RMSE comparison and the persistence bar at both horizons, was coded before the first pass. Two elements of the retention rule were formalised after the pass-1 scores had been computed: the 5\% tie band and the comparator declarations for the non-nested modules (M2 against M1, M4 against M3). Neither decides any retention outcome, the smallest pooled deficit against persistence being 17.1\%. The capelin pass is the one place where a model was added after scoring began (the index multiplier), and it is reported as a variant, not as a member of the ladder. Specification B was run last and never pooled with Specification A; the constraint is recorded in the code's docstring and in \texttt{meta.json}.

\section{Origin sets and windows}\label{si:origins}

\begin{table}[!htb]
\centering\small
\caption{Origin sets. All baselines are evaluated on the module's own origin set at each horizon. The mixed-origin persistence values that the archived Specification B summary also carries (87.6 and 317.7\,kt on 63 and 59 origins) are not comparators and are not used.}
\begin{tabularx}{\textwidth}{@{}>{\raggedright\arraybackslash}p{0.36\textwidth}Lcc@{}}
\toprule
Pass & Origins & $n$ ($h=1$) & $n$ ($h=5$) \\
\midrule
Specification A, ladder (all catch passes, survey start, two-regime $r$) & 1990--2014 ($h=1$); 1990--2010 ($h=5$); minimum eight training years & 25 & 21 \\
Specification A, capelin-index multiplier & 1991--2014; 1991--2010 (origins with a carried acoustic value) & 24 & 20 \\
Specification B, ladder and two-regime $r$ & 1965--2023; 1965--2019; minimum twelve training years & 59 & 55 \\
Specification B, capelin-index multiplier & 1988--2023; 1988--2019 & 36 & 32 \\
\bottomrule
\end{tabularx}
\end{table}

The capelin acoustic index is observed in 1982 and 1985--2023 with gaps (1983, 1984, 1993--1995, 1997, 1998, 2006, 2016, 2020--2022 missing; 2023 $= 331.3\kt$). The multiplier pass carries the last observed value across gaps and drops origins with no carried value; the two-regime pass uses the 1991 capelin collapse as a declared break and allows $r$ to differ before and after it, with the post-break value estimated only once post-break transitions are in the training window.

\FloatBarrier
\section{Uncertainty layer: full grids}\label{si:uncertainty}

Table~\ref{tab:si_dm} gives every module--benchmark pair of the archived uncertainty layer: the DM statistic (loss differences of squared errors; autocorrelation-consistent variance truncated at lag $h-1$; reported as a $z$-score), the paired moving-block bootstrap interval of the RMSE gap (block $\max(h,3)$, 20\,000 replicates, seed 0) and its two-sided $p$-value. The bootstrap in the archived script draws block indices from a generator seeded identically for the two loss series, so the resample is paired; re-running it with an explicitly shared index draw reproduces every interval to the printed digit. Table~\ref{tab:si_blocks} gives the block-length sensitivity (5\,000 replicates per cell) for the M1 and M3 persistence margins.

\begin{table}[!htb]
\centering\scriptsize
\setlength{\tabcolsep}{3.5pt}
\caption{Archived uncertainty layer: all pairs. RMSE in kt; gap $=\RMSE(A)-\RMSE(B)$; CI is the 95\% paired block-bootstrap interval of the gap; $p$ is the two-sided bootstrap $p$-value. The M2-against-M1b rows are an alternative-comparator disclosure, not part of the rule.}
\label{tab:si_dm}
\begin{tabular}{@{}ccllrrrrrlr@{}}
\toprule
Spec. & $h$ & A & B & $n$ & RMSE(A) & RMSE(B) & gap & DM $z$ & 95\% CI & $p$ \\
\midrule
A & 1 & M1 & persistence & 25 & 120.5 & 98.0 & +22.5 & 1.13 & [-13.4, +52.2] & 0.19 \\
A & 1 & M1b & persistence & 25 & 114.8 & 98.0 & +16.8 & 0.98 & [-19.4, +59.6] & 0.38 \\
A & 1 & M2 & persistence & 25 & 144.1 & 98.0 & +46.0 & 1.29 & [-12.3, +83.0] & 0.45 \\
A & 1 & M3 & persistence & 25 & 134.6 & 98.0 & +36.6 & 0.99 & [-16.8, +64.2] & 0.93 \\
A & 1 & M4 & persistence & 25 & 195.6 & 98.0 & +97.5 & 1.53 & [-9.1, +172.7] & 0.28 \\
A & 1 & M2 & M1 & 25 & 144.1 & 120.5 & +23.5 & 0.91 & [-60.7, +66.9] & 0.68 \\
A & 1 & M4 & M3 & 25 & 195.6 & 134.6 & +60.9 & 1.21 & [+4.4, +134.4] & $<0.005$ \\
A & 1 & M2 & M1b & 25 & 144.1 & 114.8 & +29.3 & 0.91 & [-67.6, +82.8] & 0.70 \\
A & 5 & M1 & persistence & 21 & 288.7 & 264.7 & +24.0 & 1.84 & [-43.1, +56.4] & 0.28 \\
A & 5 & M1b & persistence & 21 & 288.6 & 264.7 & +23.9 & 1.84 & [-43.1, +56.4] & 0.28 \\
A & 5 & M2 & persistence & 21 & 398.2 & 264.7 & +133.5 & 1.10 & [-21.0, +217.5] & 0.42 \\
A & 5 & M3 & persistence & 21 & 366.4 & 264.7 & +101.6 & 1.12 & [-0.8, +151.4] & 0.08 \\
A & 5 & M4 & persistence & 21 & 488.2 & 264.7 & +223.5 & 1.14 & [-20.8, +409.4] & 0.40 \\
A & 5 & M2 & M1 & 21 & 398.2 & 288.7 & +109.5 & 0.98 & [-77.1, +232.8] & 0.81 \\
A & 5 & M4 & M3 & 21 & 488.2 & 366.4 & +121.9 & 1.15 & [-27.4, +263.2] & 0.71 \\
A & 5 & M2 & M1b & 21 & 398.2 & 288.6 & +109.7 & 0.98 & [-77.1, +232.8] & 0.81 \\
B & 1 & M1 & persistence & 59 & 119.5 & 84.4 & +35.0 & 2.81 & [+2.7, +70.8] & 0.03 \\
B & 1 & M1b & persistence & 59 & 151.6 & 84.4 & +67.2 & 3.60 & [+18.3, +117.7] & 0.01 \\
B & 1 & M2 & persistence & 59 & 166.0 & 84.4 & +81.6 & 3.22 & [+36.4, +122.8] & $<0.005$ \\
B & 1 & M3 & persistence & 59 & 126.7 & 84.4 & +42.3 & 1.85 & [+1.0, +92.5] & 0.04 \\
B & 1 & M4 & persistence & 59 & 205.7 & 84.4 & +121.3 & 3.20 & [+53.7, +193.0] & $<0.005$ \\
B & 1 & M2 & M1 & 59 & 166.0 & 119.5 & +46.6 & 1.78 & [-20.5, +104.5] & 0.17 \\
B & 1 & M4 & M3 & 59 & 205.7 & 126.7 & +79.0 & 3.07 & [+35.6, +116.4] & $<0.005$ \\
B & 1 & M2 & M1b & 59 & 166.0 & 151.6 & +14.4 & 0.51 & [-67.5, +89.1] & 0.77 \\
B & 5 & M1 & persistence & 55 & 431.9 & 300.0 & +131.9 & 2.06 & [+33.2, +218.6] & 0.01 \\
B & 5 & M1b & persistence & 55 & 445.5 & 300.0 & +145.5 & 1.80 & [+18.2, +250.9] & 0.02 \\
B & 5 & M2 & persistence & 55 & 1058.9 & 300.0 & +758.9 & 2.15 & [+363.6, +1038.5] & $<0.005$ \\
B & 5 & M3 & persistence & 55 & 930.1 & 300.0 & +630.2 & 2.39 & [+352.8, +845.6] & $<0.005$ \\
B & 5 & M4 & persistence & 55 & 1030.7 & 300.0 & +730.8 & 2.35 & [+407.3, +978.3] & $<0.005$ \\
B & 5 & M2 & M1 & 55 & 1058.9 & 431.9 & +627.0 & 1.97 & [+195.7, +929.7] & $<0.005$ \\
B & 5 & M4 & M3 & 55 & 1030.7 & 930.1 & +100.6 & 1.88 & [+20.2, +177.4] & 0.01 \\
B & 5 & M2 & M1b & 55 & 1058.9 & 445.5 & +613.4 & 1.93 & [+170.9, +945.1] & $<0.005$ \\
\bottomrule
\end{tabular}
\end{table}

\begin{table}[!htb]
\centering\scriptsize
\setlength{\tabcolsep}{2.6pt}
\caption{Block-length sensitivity of the paired bootstrap interval (kt) for the persistence margins of M1 and M3. The archived layer uses block 3 at $h=1$ and block 5 at $h=5$. One classification changes with the block: the M3 one-step margin on Specification B includes zero for blocks of four or more.}
\label{tab:si_blocks}
\begin{tabular}{@{}cclrcccccc@{}}
\toprule
Spec. & $h$ & Module & gap & block 2 & block 3 & block 4 & block 5 & block 6 & block 8 \\
\midrule
A & 1 & M1 & +22.5 & [-13, +49] & [-13, +52] & [-14, +54] & [-16, +55] & [-18, +56] & [-18, +56] \\
A & 1 & M3 & +36.6 & [-17, +77] & [-17, +65] & [-16, +56] & [-15, +54] & [-14, +55] & [-14, +55] \\
A & 5 & M1 & +24.0 & [-37, +56] & [-37, +56] & [-41, +58] & [-43, +56] & [-51, +55] & [-65, +56] \\
A & 5 & M3 & +101.6 & [-17, +181] & [-10, +179] & [-0.4, +167] & [-0.8, +151] & [-3, +158] & [-3, +152] \\
B & 1 & M1 & +35.0 & [+5, +67] & [+2, +71] & [+1, +74] & [+1, +74] & [+0.1, +75] & [+1, +73] \\
B & 1 & M3 & +42.3 & [+4, +87] & [+1, +95] & [-1, +95] & [-3, +93] & [-4, +90] & [-4, +81] \\
B & 5 & M1 & +131.9 & [+62, +198] & [+49, +208] & [+37, +215] & [+33, +221] & [+25, +224] & [+13, +232] \\
B & 5 & M3 & +630.2 & [+436, +803] & [+394, +830] & [+362, +848] & [+348, +843] & [+310, +850] & [+272, +849] \\
\bottomrule
\end{tabular}
\end{table}

The DM $z$-statistic is less stable than the bootstrap under the truncation lag: for the M1 persistence margin at $h=1$ it moves from 1.14 at lag 0 to 2.30 at lag 3 on Specification A and from 2.81 to 1.70 on Specification B, crossing the conventional 1.96 in opposite directions, whereas the bootstrap $p$ moves within 0.17--0.19 (A) and 0.036--0.049 (B) over blocks 3--9. Four of the 28 persistence and comparator rows have a bootstrap interval excluding zero with $|z|<1.96$ (A, $h=1$, M4 vs M3; B, $h=1$, M3 vs persistence; B, $h=5$, M1b vs persistence; B, $h=5$, M4 vs M3). No verdict rests on either statistic where they disagree; the retention rule is applied to point estimates.

\FloatBarrier
\section{Period decomposition: five-year horizon, annual pass, pooled margins}\label{si:periods}

\begin{table}[!htb]
\centering\small
\caption{Five-year rolling-origin RMSE (kt) by period of the forecast origin, with origin-matched persistence; superscripts are the number of origins at which the forecast sits at the numerical floor. Companion to Table 6.}
\label{tab:si_h5}
\begin{tabular}{@{}llrrrrrrr@{}}
\toprule
Spec. & Origins & $n$ & persist & M1 & M1b & M2 & M3 & M4 \\
\midrule
A & 1990--1995 & 6 & 479.4 & 501.1$^{4}$ & 500.8$^{4}$ & 740.0$^{1}$ & 674.5$^{3}$ & 909.2$^{1}$ \\
A & 1996--2012 & 15 & 78.6 & 127.6$^{15}$ & 127.6$^{15}$ & 54.7 & 77.1 & 54.9 \\
A & 1996--2014 & 15 & 78.6 & 127.6$^{15}$ & 127.6$^{15}$ & 54.7 & 77.1 & 54.9 \\
\midrule
B & 1965--1990 & 26 & 398.8 & 573.4$^{10}$ & 595.2$^{17}$ & 1186.2$^{7}$ & 1178.6$^{7}$ & 1252.0$^{7}$ \\
B & 1991--1995 & 5 & 301.9 & 211.2$^{4}$ & 27.5$^{5}$ & 1490.1 & 834.3$^{4}$ & 1223.5$^{3}$ \\
B & 1996--2012 & 17 & 123.3 & 178.5$^{17}$ & 178.5$^{17}$ & 250.9 & 217.3 & 217.8 \\
B & 2013--2023 & 7 & 119.9 & 367.8$^{5}$ & 406.6$^{7}$ & 1358.2 & 1012.9 & 1157.9 \\
B & 1996--2023 & 24 & 122.3 & 249.0$^{22}$ & 266.1$^{24}$ & 763.3 & 576.8 & 651.7 \\
\bottomrule
\end{tabular}
\end{table}

\begin{table}[!htb]
\centering\small
\caption{Specification A, annual-landings pass: rolling-origin RMSE (kt) by period at $h=1$ (upper block) and $h=5$ (lower block). The low-biomass band gives the same ordering as the coarse-regime pass.}
\label{tab:si_annual}
\begin{tabular}{@{}llrrrrrrr@{}}
\toprule
Spec. & Origins & $n$ & persist & M1 & M1b & M2 & M3 & M4 \\
\midrule
A & 1990--1995 & 6 & 193.6 & 209.6$^{3}$ & 187.1$^{4}$ & 325.6 & 312.1$^{3}$ & 419.1 \\
A & 1996--2012 & 17 & 19.6 & 63.0$^{12}$ & 69.9$^{13}$ & 17.2 & 15.0 & 19.8 \\
A & 2013--2014 & 2 & 66.7 & 127.2 & 134.7 & 33.5 & 28.7 & 40.2 \\
A & 1996--2014 & 19 & 28.5 & 72.5$^{12}$ & 79.2$^{13}$ & 19.6 & 17.0 & 22.9 \\
\midrule
A & 1990--1995 & 6 & 479.4 & 501.1$^{4}$ & 500.8$^{4}$ & 732.3$^{1}$ & 647.5$^{3}$ & 906.6$^{1}$ \\
A & 1996--2012 & 15 & 78.6 & 127.6$^{15}$ & 127.6$^{15}$ & 50.9 & 73.0 & 49.7 \\
A & 1996--2014 & 15 & 78.6 & 127.6$^{15}$ & 127.6$^{15}$ & 50.9 & 73.0 & 49.7 \\
\bottomrule
\end{tabular}
\end{table}

\begin{table}[!htb]
\centering\small
\caption{Post-1995 pooled margins against persistence (origins 1996 onward), with paired block-bootstrap intervals (block $\max(h,3)$, 20\,000 replicates). Specification A: 1996--2014 ($h=1$) and 1996--2010 ($h=5$); Specification B: 1996--2023 and 1996--2019. The bands were chosen after the data were seen; this table describes, it does not test.}
\label{tab:si_post1995}
\begin{tabular}{@{}cccllrrrlr@{}}
\toprule
Spec. & $h$ & $n$ & Module & RMSE & persist & gap & rel. & 95\% CI (kt) & $p$ \\
\midrule
A & 1 & 19 & M1 & 73.1 & 28.5 & +44.6 & +156\% & [+28.9, +60.3] & $<0.005$ \\
A & 1 & 19 & M1b & 79.8 & 28.5 & +51.3 & +180\% & [+28.9, +69.5] & $<0.005$ \\
A & 1 & 19 & M2 & 19.4 & 28.5 & -9.1 & -32\% & [-17.1, +6.6] & 0.47 \\
A & 1 & 19 & M3 & 16.7 & 28.5 & -11.8 & -41\% & [-20.5, +3.3] & 0.19 \\
A & 1 & 19 & M4 & 22.1 & 28.5 & -6.4 & -22\% & [-14.1, +9.0] & 0.72 \\
A & 5 & 15 & M1 & 127.6 & 78.6 & +49.0 & +62\% & [+15.4, +68.1] & $<0.005$ \\
A & 5 & 15 & M1b & 127.6 & 78.6 & +49.0 & +62\% & [+15.4, +68.1] & $<0.005$ \\
A & 5 & 15 & M2 & 54.7 & 78.6 & -23.8 & -30\% & [-27.5, +11.9] & 0.76 \\
A & 5 & 15 & M3 & 77.1 & 78.6 & -1.5 & -2\% & [-2.7, +31.7] & 0.19 \\
A & 5 & 15 & M4 & 54.9 & 78.6 & -23.6 & -30\% & [-27.2, +12.1] & 0.78 \\
B & 1 & 28 & M1 & 91.5 & 41.5 & +50.0 & +121\% & [+26.3, +78.8] & $<0.005$ \\
B & 1 & 28 & M1b & 124.0 & 41.5 & +82.5 & +199\% & [+53.2, +117.7] & $<0.005$ \\
B & 1 & 28 & M2 & 102.7 & 41.5 & +61.2 & +148\% & [-2.5, +93.3] & 0.06 \\
B & 1 & 28 & M3 & 62.3 & 41.5 & +20.9 & +50\% & [-11.1, +47.0] & 0.25 \\
B & 1 & 28 & M4 & 87.4 & 41.5 & +45.9 & +111\% & [-4.9, +73.3] & 0.08 \\
B & 5 & 24 & M1 & 249.0 & 122.3 & +126.7 & +104\% & [+27.3, +196.8] & $<0.005$ \\
B & 5 & 24 & M1b & 266.1 & 122.3 & +143.8 & +118\% & [+27.3, +227.5] & $<0.005$ \\
B & 5 & 24 & M2 & 763.3 & 122.3 & +641.0 & +524\% & [+97.9, +972.2] & $<0.005$ \\
B & 5 & 24 & M3 & 576.8 & 122.3 & +454.5 & +372\% & [+46.6, +708.8] & $<0.005$ \\
B & 5 & 24 & M4 & 651.7 & 122.3 & +529.4 & +433\% & [+61.7, +813.0] & $<0.005$ \\
\bottomrule
\end{tabular}
\end{table}

Band-level intervals for the collapse bands ($n=5$--6) are degenerate under a block bootstrap and are not reported. A leave-one-origin-out check on the Specification B collapse band confirms that the orderings there do not depend on a single origin: M1 and M1b are below persistence in all five leave-one-out sets and M3 in four of five, but these are the floor-forecast wins discussed in Section 4.3.
